\documentclass[10pt,a4paper]{amsart}
\usepackage[margin=19mm]{geometry}
\usepackage{amsmath,amssymb,mathtools}
\usepackage[scr=rsfs,cal=euler]{mathalfa}
\usepackage{xfrac}
\usepackage[unicode,hidelinks,bookmarks=false]{hyperref}
\newcommand{\ie}{i.e.\ }
\newcommand{\cf}{cf.\ }
\newcommand{\resp}{respectively\ }
\newcommand{\Subsection}{\subsection}
\providecommand{\nobreakdash}{\nobreak\hbox{-}}
\setlength{\emergencystretch}{2em}
\multlinegap0pt
\usepackage{bm}
\usepackage{bbm}
\usepackage{dsfont}
\usepackage{xspace}
\usepackage{cancel}
\usepackage{mathtools}
\usepackage{amsthm}
\makeatletter
\@ifundefined{theorem}{\newtheorem{theorem}{Theorem}}{}
\makeatother
\title[Exact solution of the (2+1)-dimensional damping forcing coup]{Exact solution of the (2+1)-dimensional damping forcing coupled Burgers equation by using Darboux transformation}
\author{Prasanta Chatterjee, Nanda Kanan Pal, Dipan Saha, Santanu Raut}
\date{Source: arXiv:2601.13634v1 (CC BY 4.0)}
\begin{document}
\maketitle

\noindent\emph{Source and licence.} Exact solution of the (2+1)-dimensional damping forcing coupled Burgers equation by using Darboux transformation. Version arXiv:2601.13634v1 (CC BY 4.0), distributed under the Creative Commons Attribution 4.0 International licence, as recorded in the arXiv licence metadata for this exact version and shown on the abstract page https://arxiv.org/abs/2601.13634v1. Full rights notice: licenses/arxiv-2601.13634-CC-BY-4.0.txt.\par\smallskip

\noindent\textit{Editorial scope.} Counted unit: the Theorem whose source label is \texttt{Nsol} in the article (source0.tex lines 177--192, source label \texttt{Nsol}), delivered together with the article's own complete original proof of it (source0.tex lines 193--275, one \texttt{proof} environment of 83 lines). No mathematics is written, repaired or re-derived: statement and proof are the article's own text line for line. Required context retained uncounted: laxpair (lines 93--95); DT (lines 120--122); nlaxpair (lines 170--172). Every definition, display and label of those lines is retained unchanged. Omitted from this package and not redistributed: 1--92, 96--119, 123--169, 173--176, 276--376. Nothing inside a retained excerpt is omitted or altered: every retained body is delivered exactly as the article's lines are written, so no omission receipt of any kind accompanies this package. Citations: the retained excerpts carry no citation command at all, so no bibliography entry of the article is transcribed and no reference is redistributed. Reference adaptation: none was needed and none was made; every reference inside the delivered unit resolves inside it, so no unresolved reference or citation is produced. Printed slips kept literal: no printed word, symbol, hypothesis or number of the counted statement or of its proof was altered. Layout and class: the article's own class is \texttt{article} with packages including authblk, amsmath, amsfonts, amssymb, bm, geometry, scalefnt, xcolor, hyperref, graphicx, epstopdf, subfigure, fullpage, float; the wrapper is the lane's pinned \texttt{amsart} editorial wrapper at 10pt on A4, which restates the theorem-like environments the retained text needs and adds the article's own macro definitions that the retained text needs (none), with extra wrapper packages (bm, bbm, dsfont, xspace, cancel, mathtools). The delivered numbering is the unit's own. Assets: no retained span contains a figure environment, a graphics-inclusion command, a PDF-page inclusion, a table or a TikZ picture; no image file is copied into this package, the package declares no asset dependency and no image file is redistributed with it. Licence: version arXiv:2601.13634v1 of this article is distributed under the Creative Commons Attribution 4.0 International licence, as recorded in the arXiv licence metadata for this exact version and shown on the abstract page https://arxiv.org/abs/2601.13634v1. A full rights notice is delivered at licenses/arxiv-2601.13634-CC-BY-4.0.txt. Characters: every retained excerpt is pure ASCII, so the delivered engine, XeLaTeX with its default Latin Modern text font, reports no missing character.
\begin{equation}
L\psi=\psi_y,\hspace{5mm}\psi_t=M\psi,\label{laxpair}
\end{equation}

\begin{equation}
	\psi[1]=(\partial_x-\partial_x\log\psi_1)\psi=\psi_{x}+B\psi, \label{DT}
\end{equation}

\begin{equation}
 L[1]\psi[1]=\psi_y[1],\hspace{2mm}\psi_t[1]=M[1]\psi[1]. \label{nlaxpair}
\end{equation}

\begin{theorem}
	Let $\psi_{1}$, $\psi_{2}$, $\psi_{3}$, $\cdots$, $\psi_{N}$ be the non-zero particular solutions of the linear equations \eqref{laxpair}. Then the new eigen function  $\psi[N]=\frac{W(\psi_{1},\psi_{2},\psi_{3},\cdots,\psi_{N},\psi)}{W(\psi_{1},\psi_{2},\psi_{3},\cdots,\psi_{N})}$ satisfy the Lax pair  $L[N]\psi[N]=\psi_y[N] $,\quad $\psi_t[N]=M[N]\psi[N]$ corresponding to the potential stated below
\begin{equation}
\begin{split}
&u[N]=u+3\Lambda(t)\partial^{2}_{x}\log W(\psi_{1},\psi_{2},\psi_{3},\cdots,\psi_{N}),\\
&v[N]=v+Nu_x+3\Lambda(t)\sum_{r=1}^{N}\bigg(\partial_x\log \frac{W(\psi_1,\psi_2,\cdots,\psi_r)}{W(\psi_1,\psi_2,\cdots,\psi_{r-1})}\partial^{2}_x\log\frac{W(\psi_1,\psi_2,\cdots,\psi_r)}{W(\psi_1,\psi_2,\cdots,\psi_{r-1})}+\partial^{3}_x\log W(\psi_1,\psi_2,\cdots,\psi_r)\bigg)\label{Nsol}
\end{split}
\end{equation}
where
\begin{equation}
\begin{split}
&L[N]=\partial_x^3+\frac{u[N]-H(t)}{\Lambda(t)}\partial_x+\frac{v[N]-H(t)}{\Lambda(t)}I,~ M[N]=\partial_x^3+\frac{3}{2}\partial^2_x+\frac{u[N]-H(t)}{\Lambda(t)}\partial_x+\partial_y+\frac{u[N]+v[N]-2H(t)}{\Lambda(t)}I.
\end{split}
\end{equation}
and 	$W(\psi_{1},\psi_{2},\psi_{3},\cdots,\psi_{K})=\text{det}(\theta_1, \theta_2, \theta_3,\cdots, \theta_K)$, $\theta_j=(\psi_j,\partial_{x}\psi_j,\partial^2_{x}\psi_j,\cdots,\partial^{K-1}_{x}\psi_j)^T$ is the column vector forming by $\psi_j$ and its partial derivatives upto $(K-1)$ th order with respect to the variable $x$ , $j=1,2,3,\cdots, N$
\end{theorem}

\begin{proof}
We use the induction theory to prove the above result for $N$-fold Darboux transformation. For $N=1$, we proved in theorem 1, the result is true for one-fold Darboux transformation. Now we shall prove the result is true for $N=2$. Again we apply the one-fold Darboux transformation in equation \eqref{nlaxpair} defined as
\begin{equation}
\psi[2]=(\partial_x-\partial_x\log\psi_1[1])\psi[1] \label{2DT}
\end{equation}
Where $\psi_1[1]$, is the nonzero particular solution of \eqref{nlaxpair}. The value of $\psi_1[1]$, one can obtained by \eqref{DT} is $\psi_1[1]=\frac{W(\psi_{1},\psi_{2})}{\psi_{1}}$. Here $\psi_{2}$ is the non zero particular solution of \eqref{laxpair}. By equation \eqref{DT}, the expression of $\psi[2]$ in equation \eqref{2DT} can be written as in the following form
\begin{equation}
\psi[2]=\frac{W(\psi_{1},\psi_{2},\psi)}{W(\psi_{1},\psi_{2})},
\end{equation}
which is called 2-fold Darboux transformation on $\psi$ in \eqref{laxpair}. Then the new potentials becomes
\begin{equation}
\begin{split}
&u[2]=u[1]+3\Lambda(t)\partial^{2}_x\log\psi_1[1],~
v[2]=v[1]+u_x[1]+3\Lambda(t)(\partial_x\log\psi_1[1]\partial^{2}_x\log\psi_1[1]+\partial^{3}_x\log\psi_1[1]).
\end{split}
\end{equation}
Subtitutiong the values of $u[1]$, $v[1]$, and  $\psi_1[1]$ into above expression, we have
\begin{equation}
\begin{split}
&u[2]=u+3\Lambda(t)\partial^{2}_{x}\log W(\psi_{1},\psi_{2}),\\
&v[2]=v+2u_x+3\Lambda(t)\sum_{r=1}^{2}\bigg(\partial_x\log \frac{W(\psi_1,\psi_2,\cdots,\psi_r)}{W(\psi_1,\psi_2,\cdots,\psi_{r-1})}\partial^{2}_x\log\frac{W(\psi_1,\psi_2,\cdots,\psi_r)}{W(\psi_1,\psi_2,\cdots,\psi_{r-1})}+\partial^{3}_x\log W(\psi_1,\psi_2,\cdots,\psi_r)\bigg)
\end{split}
\end{equation}
Thus the result is true for $N=2$.\\
 We proved the Darboux transformation result is true for $N=1$ and $N=2$, we shall assume the result is true for $N=K$, therefore for $K$- fold Darboux transformation on $\psi$ in \eqref{laxpair} is	
\begin{equation}
\psi[K]=\frac{W(\psi_{1},\psi_{2},\psi_{3},...\psi_{K},\psi)}{W(\psi_{1},\psi_{2},\psi_{3},...\psi_{K})} \label{KDT}
\end{equation}
Thus $\psi[K]$ obey the followings equations
\begin{equation}
\begin{split}
 L[K]\psi[K]=\psi_y[K],\hspace{2mm}\psi_t[K]=M[K]\psi[K] 
 \label{covlaxpairNN}
\end{split}
\end{equation}
where the new potential becomes after $K$-fold Darboux transformation
\begin{equation}
 \begin{split}
 &u[K]=u+3\Lambda(t)\partial^{2}_{x}\log W(\psi_{1},\psi_{2},\psi_{3},\cdots,\psi_{K}),\\
 &v[K]=v+Ku_x+3\Lambda(t)\sum_{r=1}^{K}\bigg(\partial_x\log \frac{W(\psi_1,\psi_2,\cdots,\psi_r)}{W(\psi_1,\psi_2,\cdots,\psi_{r-1})}\partial^{2}_x\log\frac{W(\psi_1,\psi_2,\cdots,\psi_r)}{W(\psi_1,\psi_2,\cdots,\psi_{r-1})}+\partial^{3}_x\log W(\psi_1,\psi_2,\cdots,\psi_r)\bigg)
 \end{split}                              \label{newsolNN}
\end{equation}
and $W(\psi_{1},\psi_{2},\psi_{3},\cdots,\psi_{K})=\text{det}(\theta_1, \theta_2, \theta_3,\cdots, \theta_K)$, $\theta_j=(\psi_j,\partial_{x}\psi_j,\partial^2_{x}\psi_j,\cdots,\partial^{K-1}_{x}\psi_j)^T$ is the column vector forming by $\psi_j$ and its partial derivatives upto $(K-1)$ th order with respect to the variable $x$ , $j=1,2,3,\cdots, N$ and $\psi_{1}$, $\psi_{2}$, $\psi_{3}$, $\cdots$, $\psi_{K}$ are the nonzero solutions of \eqref{laxpair}.\\
The one-fold Darboux transformation on $\psi[K]$ in \eqref{covlaxpairNN} defined as
\begin{equation}
\psi[K+1]=(\partial_x-\partial_x\log\psi_K[K])\psi[K],\label{DTK}
\end{equation}
where $\psi_K[K]$ is the non zero particular solution of \eqref{covlaxpairNN}. Thus the value of $\psi_K[K]$ one can obtained through \eqref{KDT} as
\begin{equation}
\psi_K[K]=\frac{W(\psi_{1},\psi_{2},\psi_{3},...,\psi_{K},\psi_{K+1})}{W(\psi_{1},\psi_{2},\psi_{3},...,\psi_{K})}
\end{equation}
where $\Psi_{K+1}$ be the nonzero particular solution of \eqref{laxpair}. Substituing the values of  $\psi[K]$ and $\psi_K[K]$  into \eqref{DTK}, thus we obtained $(K+1)$-fold Darboux transformation on $\Psi$ as 
\begin{equation}
\Psi[K+1]=\frac{W(\psi_{1},\psi_{2},\psi_{3},...,\psi_{K+1},\psi)}{W(\psi_{1},\psi_{2},\psi_{3},...\psi_{K+1})} \label{MKDT}
\end{equation}
Therefore $\psi[K+1]$ obey the following equations
\begin{equation}
\begin{split}
L[K+1]\psi[K+1]=\psi_y[K+1],\hspace{2mm}\psi_t[K+1]=M[K+1]\psi[K+1] 
 \label{Mcovlaxpair}
\end{split}
\end{equation}
thus the new potential becomes after (K+1)-fold Darboux transformation as
\begin{equation}
\begin{split}
u[K+1]=&u[K]+3\Lambda(t)\partial^{2}_{x}\log (\psi_K[K]),  \\
=&u+3\Lambda(t)\partial^{2}_{x}\log W(\psi_{1},\psi_{2},\psi_{3},\cdots,\psi_{K})\\&+3\Lambda(t)\partial^{2}_{x}\log\bigg(\frac{W(\psi_{1},\psi_{2},\psi_{3},...\psi_{K},\psi_{K+1})}{W(\psi_{1},\psi_{2},\psi_{3},...\psi_{K})}\bigg),\\
=&u+3\Lambda(t)\partial^{2}_{x}\log W(\psi_{1},\psi_{2},\psi_{3},\cdots,\psi_{K+1}).
\end{split}                                             \label{Knewsol}
\end{equation}
and 
\begin{equation}
\begin{split}
&v[K+1]=v[K]+u_x[K]+ 3\Lambda(t)(\partial_x\log \psi_K[K]\partial^{2}_x\log \psi_K[K]+\partial^{3}_x\log \psi_K[K])
\end{split}
\end{equation}
By using the values of $v[K]$, $u[K]$, and $\psi_K[K]$, the above expression reduced to 
\begin{equation}
v[K+1]=v+(K+1)u_x+3\Lambda(t)\sum_{r=1}^{K+1}\bigg(\partial_x\log \frac{W(\psi_1,\psi_2,\cdots,\psi_r)}{W(\psi_1,\psi_2,\cdots,\psi_{r-1})}\partial^{2}_x\log\frac{W(\psi_1,\psi_2,\cdots,\psi_r)}{W(\psi_1,\psi_2,\cdots,\psi_{r-1})}+\partial^{3}_x\log W(\psi_1,\psi_2,\cdots,\psi_r)\bigg)
\end{equation}
Thus the result is true for $N=K+1$.\\
This complete the proof.
\end{proof}

\end{document}
